\startreportpage{25.4}
\templaterow{结\quad 论\\（结\quad 果）}{94}{
UDP实测请求首部8字节，Length为39，等于8字节首部加31字节数据；外层IPv4总长59，等于20+39。源/目的端口、长度、校验和的原始字节与Wireshark解释一致。响应端点反向，UDP Length增为85，数据为77字节。请求与响应相隔9.014毫秒。请求校验和验证正确，响应零校验和在IPv4中表示未使用。

同一条TCP连接中，客户端主动三次握手建立连接，两端SYN各消耗一个序号。客户端/服务器MSS分别1460/1300，Window Scale分别为左移6/7位；SYN报文本身的窗口不缩放，后续报文按协商因子解释。连续上传帧5—7的数据长度均1288，Seq逐段递增；帧14的Ack 4543累计确认至序号4542，对帧7的RTT为269.348毫秒。

POST携带公开文件\code{alice.txt}，重组文件152138字节与原文件完全一致，服务器返回HTTP 200及成功页面。服务器首先发送FIN，后续响应/重传补齐抓包中的序号缺口，客户端累计确认后再发FIN，最终收到服务器ACK。传输虽有重复ACK及重传，仍完成应用交付与正常关闭。
}
\templaterow{小\quad 结}{83}{
本次实验把字段值、原始字节、长度计算和应用结果联系起来。UDP Length与帧长属于不同层次；端口是16位字段，理论长度上限也要说明约束条件。伪首部参与校验计算，但不是在线上另加12字节首部；Unverified与校验错误也应区别处理。

TCP分析中，MSS、实际数据长度、原始窗口与缩放窗口不能互相替代。ACK指向下一个期望序号，既可能累计确认多个数据段，也可能在有缺口时保持不变；本次SACK与随后Ack 782清楚体现了这一点。关闭分析应保持原始时间顺序，并结合序号范围，不应把实际报文改写成理想示意图。

本次证据仅来自本机\code{en0}，不能据此定位丢包发生在哪一跳，也不能由抓包工具零丢弃断言线路无丢包。实验未调整系统DNS或代理；直接HTTP浏览器使用独立临时配置。附件保存最终pcapng、成功采集原文件、筛选映射、命令输出、字段详情和完整截图，便于复核；正文的截图裁剪坐标另有记录。
}
\TemplateReview
